\scanpage{349}
\mspage{10}
These matrices form a nilpotent Lie group structure \(H\) on \(R^3\) whose invariant forms are generated by say
\[
 dx,\ dy,\ dz+xdy.
\]
The differential algebra generated by these is isomorphic to the model \(\Lambda(x_1,x_2,y)\).

By integrating along non-closed paths starting at \(p\) we obtain a map \(\phi=\phi(\omega_1,\omega_2,\eta)\) of the manifold \(M\) into the compact 3-dimensional nil-manifold \(H/\Gamma\),
\diagram{\begin{tikzcd}[ampersand replacement=\&]
 M\handwritten{\arrow[r]}\arrow[r,phantom,"\phi"]\& H/\Gamma
\end{tikzcd}}
where \(\Gamma\) is \((\pi_1/[[\pi_1,\pi_1],\pi_1])/\text{torsion}\). In this case \(H/\Gamma\) would be
\[
 \mathopen{\handwritten{\Bigg(}}\begin{matrix}1&x&z\\0&1&y\\0&0&0\end{matrix}\mathclose{\handwritten{\Bigg)}}
 \handwritten{\Bigg/}
 \mathopen{\handwritten{\Bigg(}}\begin{matrix}1&n&\ell/k\\0&1&m\\0&0&0\end{matrix}\mathclose{\handwritten{\Bigg)}}
\]
where \(x,y,z\) are real numbers and \(n,m,\ell,k\) are integers with \(k\) fixed.

The composition
\diagram{\begin{tikzcd}[ampersand replacement=\&]
 M\arrow[r,"\phi"] \& H/\Gamma\handwritten{\arrow[r]}\arrow[r,phantom,"\text{abel.}"] \& \text{Torus}
\end{tikzcd}}
would be the Jacobian determined by \(\omega_1\) and \(\omega_2\).

These non-abelian periods generalize to the higher nilpotent quotents \(\{\Gamma_k\}\) of \(\pi_1\) and analogous nilpotent quotients of the higher homotopy groups. (Regarded as \(\pi_1\)-modules). For example in the simply connected case we obtain a complete set of \uline{real} periods for the higher homotopy groups.
% archive sheet number: 350
